\begin{definition}[Inf-convolution, (3.13)]
Let $\mathbb G$ be a Carnot setting (\noderef{CarnotSetting}) with group law $*$ and
homogeneous norm $\|\cdot\|$, put $q:=2\,r!$, let $\Omega\subset\mathbb R^n$, let
$w:\mathbb R^n\to\mathbb R$, let $\varepsilon>0$ and let $x\in\mathbb R^n$. The
\emph{inf-convolution} of $w$ is
\[
 w_{\varepsilon}(x)\ :=\ \inf\Big\{\,w(z*x)+\frac{\|z\|^{q}}{2\varepsilon}\ :\
   z\in\mathbb R^n,\ z*x\in\overline\Omega\,\Big\} .
\]

This is the inf-convolution of (3.13),
$w_{\varepsilon}(x)=\inf_{y\in\overline\Omega}\{w(y)+\|y*x^{-1}\|^{q}/2\varepsilon\}$,
after the substitution $y=z*x$, so that $w_{\varepsilon}$ is an infimum of \emph{left}
translates of $w$; it is the form in which the paper uses it in the proof of Lemma
Lemma~3.8(4). One has $w_{\varepsilon}=-\big((-w)^{\varepsilon}\big)$, where
$(-w)^{\varepsilon}$ is the sup-convolution of $-w$, which is why the inf-convolution of a
super-solution is semi-concave where the sup-convolution of a sub-solution is
semi-convex.

The infimum is taken over the set of admissible $z$, which is non-empty exactly when
$x\in\overline\Omega$, since $z=0$ gives $0*x=x$ and $\|0\|=0$ --- in which case
$w_{\varepsilon}(x)\le w(x)$. It is finite whenever $\Omega$ is bounded and $w$ is
continuous on $\overline\Omega$, because the family is then bounded below by
$\inf_{\overline\Omega}w$. In the formal statement the infimum is the infimum of the
indexed family of real numbers, which by convention returns $0$ when the family is empty or
unbounded below; the nodes that use $w_{\varepsilon}$ carry the hypotheses $\Omega$
bounded, $w\in C(\overline\Omega)$ and $x\in\overline\Omega$ under which neither
degenerate case occurs.
\end{definition}
